% ============================================================
%  第3章  硬件设计与实现
% ============================================================
\chapter{硬件设计与实现}

本章详细阐述头戴式眼部康复装置的硬件系统设计。装置以DM-MC02核心板为主控平台，围绕STM32H723VGT6高性能微控制器构建外围电路，集成了惯性姿态感知、双轴舵机云台、激光视觉刺激、语音交互、听觉报警和患者主动反馈等硬件模块，最终在自行车头盔上完成系统集成。

\section{主控电路与最小系统}

\subsection{DM-MC02核心板}

系统的开发平台选取了深圳达妙科技有限公司出品的DM-MC02电机控制核心板~\cite{dmmc02_um}。该板以STMicroelectronics公司的STM32H723VGT6微控制器为核心~\cite{stm32h7_rm,stm32h7_ds}，最高可运行在550MHz的时钟频率下，片内固化了1MB容量的Flash存储器和564KB的SRAM，并且同时配备有单精度和双精度两套FPU浮点运算硬件单元，在惯性传感器信号融合与实时控制类算法的处理上展现出了充足的算力裕量。

核心板外形尺寸约为67mm$\times$44mm，净重约25g，较好地满足了头部可穿戴类设备对于体积和质量的双重限制。板上集成了一颗24MHz高频石晶振荡器（HSE）和一颗32.768kHz低频晶体（LSE），分别承担系统主时钟生成和RTC计时时钟供给的功能；此外还配置了一片Flash存储芯片和一颗WS2812规格的可编程RGB LED，前者为固件和数据的本地存储提供了充足容量，后者则为显示系统的运行状态开辟了通道。DM-MC02开发板的三维渲染图参见图\ref{fig:dm_mc02}。

\begin{figure}[h!]
\centering
\includegraphics[width=0.9\textwidth]{dm-mc02_实物.pdf}
\caption{DM-MC02开发板3D渲染图}
\label{fig:dm_mc02}
\end{figure}

\subsection{系统时钟配置}

系统的时钟树从外部24MHz HSE晶振出发~\cite{stm32h7_rm}，经PLL倍频锁相后进行分配。PLL设置为：M分频系数为2、N倍频系数为40、P分频系数为1，产生480MHz的PLL输出频率，作为系统主时钟（SYSCLK=480MHz）的来源。AHB总线经HPRE=2分频后工作于240MHz；APB1和APB2总线经各2分频后工作于120MHz。TIM1定时器挂载于APB2总线，受益于STM32H7的定时器时钟倍频机制，TIM1的实际计数时钟达到240MHz。各外设的时钟频率均满足或远超其功能需求。

\subsection{系统引脚配置}

根据DM-MC02核心板所引出的排针阵列及以上各功能模块的实际接线需求，笔者在设计过程中详细考虑了引脚分配问题，表\ref{tab:pinout}汇总给出了系统主要引脚的分配结果，图\ref{fig:dm-mc02}则展示了该开发板的管脚标注。%引脚规划过程中遵循了几条基本原则：各外设尽可能占用其默认功能的复用引脚以降低软件初始化阶段的配置复杂程度；属于同一功能组的控制信号尽量集中安排在同一个GPIO端口内以缩短跨端口的读写开销；此外还为后续的功能扩展预留了一定数量的空闲引脚。

\begin{table}[h!]
\centering
\caption{系统主要引脚配置表}
\label{tab:pinout}
\begin{tabular}{p{2.2cm} p{4cm} p{2.8cm} p{4cm}}
\toprule[1.5pt]
\textbf{引脚} & \textbf{外设功能} & \textbf{连接模块} & \textbf{说明} \\
\midrule[0.75pt]
PC0    & ACC\_CS                & BMI088加速度计 & 加速度计片选，低电平有效 \\
PC1    & SPI2\_MOSI             & BMI088（共用）  & 主机输出/从机输入数据线 \\
PC2    & SPI2\_MISO             & BMI088（共用）  & 主机输入/从机输出数据线 \\
PC3    & GYRO\_CS               & BMI088陀螺仪   & 陀螺仪片选，低电平有效 \\
PB13   & SPI2\_SCK              & BMI088（共用）  & 串行时钟，7.5\,Mbps \\
PE10   & GPIO\_EXTI10           & BMI088加速度计 & 加速度计数据就绪中断 \\
PE12   & GPIO\_EXTI12           & BMI088陀螺仪   & 陀螺仪数据就绪中断 \\
PE9    & TIM1\_CH1              & Y轴舵机 & 50\,Hz PWM \\
PE13   & TIM1\_CH3              & X轴舵机 & 50\,Hz PWM \\
PA0    & GPIO\_Output           & 激光模组        & 推挽输出，高电平点亮 \\
PA7    & SPI6\_MOSI             & WS2812状态灯    & SPI模拟单总线协议 \\
PE7    & UART7\_RX              & CI1302模块      & MCU接收，模块发送\\
PE8    & UART7\_TX              & CI1302模块      & MCU发送，模块接收\\
PA2    & GPIO\_Input    & 患者反馈按键    & 按下接地 \\
PB15   & TIM12\_CH2             & 板载蜂鸣器      & 控制音量 \\
PC13~15   & GPIO\_Output           & 5V电源通道     & 高电平使能 \\
%PC14   & GPIO\_Output           & 5V电源通道2     & 高电平使能 \\
%PC15   & GPIO\_Output           & 5V电源通道3     & 高电平使能 \\
PA13   & SWDIO                  & SWD调试接口     & 串行线数据 \\
PA14   & SWCLK                  & SWD调试接口     & 串行线时钟 \\
PH0/PH1 & OSC\_IN/OSC\_OUT       & 24\,MHz HSE晶振 & 外部高速时钟输入/输出 \\
\bottomrule[0.75pt]
\end{tabular}
\end{table}

\begin{figure}[h!]
\centering
\includegraphics[angle=90,width=1.0\textwidth]{dm-mc02.pdf}
\caption{开发板管脚标注图}
\label{fig:dm-mc02}
\end{figure}


\section{姿态传感器模块}

\subsection{BMI088惯性传感器}

系统的姿态感知模块采用了Bosch Sensortec公司推出的BMI088六轴惯性测量单元~\cite{bmi088_ds}。该器件内部封装了一颗16位分辨力的三轴加速度计和一颗16位分辨力的三轴陀螺仪，两颗敏感芯体在物理结构和电气接口层面保持彼此独立，各自挂载在一条专用的SPI总线上并拥有单独的片选控制引脚。这种双芯独立架构带来的实际好处在于：加速度计和陀螺仪的量程、输出数据率等参数可以分开独立设定，且两条SPI通信链路之间不存在互锁阻塞——主控芯片能够在同一时刻触发加速度计与陀螺仪的DMA读取传输，从而在相当程度上提升了惯性数据的整体刷新频率。

加速度计一侧的量程参数设置为$\pm$3g，折算后的灵敏度约为$0.897\times10^{-3}$g/LSB，输出数据率（ODR）固定在800Hz；陀螺仪一侧的量程参数设置为$\pm$2000dps，相应的灵敏度约为$1.065\times10^{-3}$ dps/LSB，数据带宽则选取了1000Hz/116Hz的组合配置。在此参数设定下，系统能够较为从容地捕获佩戴对象头部的微小角度运动以及缓慢的姿态漂移信号。
\subsection{SPI通信接口}

加速度计和陀螺仪共用SPI2总线与MCU通信，且片选信号各自独立——加速度计由PC0引脚控制，陀螺仪由PC3引脚控制。两个传感器的数据就绪中断信号分别连接PE10（ACC\_INT）和PE12（GYRO\_INT），并通过这两个引脚来通过外部中断方式通知MCU传感器数据已更新。

SPI2外设配置为全双工主模式、8位数据帧、预分频系数为32（对应SCK频率为240MHz/32=7.5Mbps）。为加速数据传输，SPI2的收发通道分别配置了DMA1\_Stream0（SPI2\_RX，外设到存储器）和DMA1\_Stream1（SPI2\_TX，存储器到外设），DMA中断优先级分别设为1和2。

BMI088惯性传感器示意图如图\ref{fig:bmi088}所示。

\begin{figure}[h!]
\centering
\includegraphics[width=0.9\textwidth]{bmi088.jpg}
\caption{BMI088惯性传感器3D渲染图及实物图}
\label{fig:bmi088}
\end{figure}

\section{舵机云台控制电路}

\subsection{SG90双轴舵机云台}

视觉引导装置的二维平面定位由双轴舵机云台实现，云台由两枚SG90型微型舵机正交安装构成，如图\ref{fig:sg90}所示。SG90舵机的工作电压为4.8V至6.0V~\cite{sg90_ds}，在4.8V供电下额定扭矩为1.2 kg·cm，转速约为60°/0.12s；重量仅为9g，适合头戴式应用对载荷的要求。舵机的转角范围为0°至180°，由标准PWM信号控制——脉冲宽度从0.5ms（对应0°）至2.5ms（对应180°），1.5ms脉宽对应中间位置90°。

X轴舵机（俯仰方向）负责激光点在垂直方向上的运动，对应患者的上下视角；Y轴舵机（偏航方向）负责激光点在水平方向上的运动，对应患者的左右视角。两个舵机的机械结构通过3D打印转接件正交连接，激光模组固定在Y轴舵机输出臂末端。
\begin{figure}[h!]
\centering
\includegraphics[width=0.9\textwidth]{servo.jpg}
\caption{SG90双轴舵机云台实物及3D模型图}
\label{fig:sg90}
\end{figure}

\subsection{PWM驱动信号生成}

SG90舵机所需的标准PWM控制信号由TIM1高级定时器负责生成，该定时器的工作时钟频率为240MHz。PWM参数的设定方案如下：将预分频器数值取为240，此举使计数器的节拍频率下降至1MHz，即每个计数脉冲对应1$\mu$s的时基；自动重装载寄存器（ARR）的取值则设定为19999，这意味着计数器每完成20000次翻转对应一个完整的PWM周期，换算后即为20ms的周期时长（对应50Hz的刷新频率），完全吻合舵机对50Hz控制信号的规格要求。脉宽的控制通过比较寄存器CCR来实现：当CCR设为500时输出0.5ms的脉宽，对应舵机的0°极限位置；CCR取值2500时输出2.5ms的脉宽，对应180°的另一个极限位置；CCR取值1500时输出1.5ms脉宽，对应90°的中间参考角度。

TIM1的第一通道（CH1）经由PE9引脚引出，用于驱动Y轴（偏航方向）的舵机；第三通道（CH3）经由PE13引脚引出，用于驱动X轴（俯仰方向）的舵机。两路PWM输出信号在初始化流程完成之后即刻被启动，以此确保舵机在上电的最初瞬间即进入锁定状态，避免因初始脉宽的不确定性而引发舵机随机抖动的现象。
\section{视觉反馈模块}

\subsection{激光指示器}

视觉反馈链路中的核心部件是一枚波长为650nm的红色激光二极管，其外围配合了一副通过3D打印工艺自制的外壳加以固定和保护。该激光模组的标准工作电压为5V，内部已经预先集成了一片小型驱动电路，仅需对其信号引脚（S）施加逻辑电平即可实现对激光管亮灭状态的便捷控制。考虑到STM32主控芯片的通用GPIO输出标称电压仅为3.3V，其驱动能力不足以直连5V电压等级的激光模组，系统在此处采用了典型的三线制接线方案：将激光模组的VCC引脚接入核心板提供的5V动力电源轨，将其GND引脚接入系统公共地网络，而信号引脚（S）则与STM32的PA0通用输出端口相连。当PA0口线置为逻辑高状态（约3.3V）时，模组内部的驱动电路自行完成5V电压与激光管的通路切换工作，使得激光管按照额定功率正常发光。

PA0引脚被设定为推挽输出工作模式，上电后的默认初始电平为逻辑低（对应激光管的熄灭状态）。上层的控制软件通过调用激光开关控制函数来分别设置和清除PA0的电平输出，从而以毫秒级的时间精度管理激光管的亮灭状态。激光模组的物理安装位置位于Y轴舵机的输出臂末端，借助双轴舵机云台在俯仰和偏航两个自由度上的协同运动，激光光斑能够在患者前方墙面的任意二维坐标处实现精确定位投影。激光模块的外形结构见图\ref{fig:laser}。
\begin{figure}[h!]
\centering
\includegraphics[width=0.9\textwidth]{LASER.jpg}
\caption{激光指示模组3D模型及实物图}
\label{fig:laser}
\end{figure}

\subsection{WS2812可编程状态指示灯}

DM-MC02核心板上集成了一枚WS2812-2020型可编程RGB LED灯珠~\cite{ws2812_ds}（其封装尺寸为2.0mm$\times$2.0mm），这枚灯珠的信号输入线连接在PA7引脚上，实物如图\ref{fig:led}所示。WS2812所采用的通信协议为单总线归零码机制，其数据传输速率约为800kHz，每一位数据由高电平脉冲的持续时间宽度来区分逻辑"0"码和逻辑"1"码。

系统将SPI6外设所引出的MOSI信号线（即PA7引脚）复用为WS2812的数据通道。LED颜色的刷新流程可以概括为：首先构建一组长度为24字节的颜色数据帧（其中绿色分量占8字节、红色分量占8字节、蓝色分量占8字节，每一比特位分别以$0x78$代表逻辑"1"、以$0x60$代表逻辑"0"进行编码），随后通过SPI6外设将这24字节的颜色数据帧完整发送出去，紧接着再补发100字节的全零数据充当复位脉冲（其时长远超50$\mu$s），WS2812芯片即可自动锁存最新收到的RGB数值并据此驱动LED显示出相应的目标颜色。

在系统全工作周期的不同阶段，WS2812的状态提示策略如下：当系统已经完成就绪但尚未进入任何训练模式时，LED被设定为绿色呼吸灯效果（其亮度按照正弦函数的映射规律进行渐变，每轮完整的呼吸循环大约持续2秒）；一旦监测到佩戴者头部姿态超出了预设的安全界限，LED的工作模式立刻从呼吸效果切换为红色常亮状态，以此向患者传递不间断的视觉警示信号；而当系统再度回到或正处于空闲就绪状态时，LED则始终维持其默认的绿色呼吸灯显示状态。

\begin{figure}[h!]
\centering
\includegraphics[width=0.95\textwidth]{LED.jpg}
\caption{WS2812 可编程状态指示灯实物图}
\label{fig:led}
\end{figure}
%\subsection{指示LED}
%为提供明确的训练模式和系统状态指示，系统外扩了两路独立LED灯珠，分别由PB8和PB9 GPIO输出直接驱动。聚焦训练LED（PB8）用于视觉聚焦模式中作为近距注视目标，同时在其他模式的空闲状态下以200ms周期闪烁表示系统就绪；状态指示LED（PB9）用于训练反馈阶段的闪烁提示。两路LED均由GPIO输出直接通过限流电阻连接至LED正极，LED负极接地，无需额外的驱动芯片。

\section{语音交互模块}

\subsection{CI1302语音识别与播报模块}

系统的语音功能部分选取了深圳亚博智能科技有限公司所生产的AI语音交互集成模块，借助该模块来完成语音指令的识别输入以及文本到语音的合成播报输出两项基本任务。这块模块的核心处理芯片是云知声公司开发的CI1302神经网络处理器，其内部固化了离线语音识别引擎和TTS合成引擎，能够处理超过110条预设在固件中的语音指令，在距离模块约5米的半径范围内可实现较高准确度的中英文指令识别。模块同时还具备回声消除和自然对话等辅助功能，出厂固件中将唤醒词预设为"你好小盈"。

语音模块同主控之间的数据交换通过UART异步串行通信接口来完成：模块一侧的UART\_RX接收线连接到主控芯片PE8引脚（UART7\_TX发送线），而模块的UART\_TX发送线则连接到主控芯片的PE7引脚（UART7\_RX接收线）。主控一侧将UART7外设的工作模式设定为异步全双工，通信波特率取值为115200bps，数据帧结构为8位数据位、1位停止位、无校验位、无硬件流量控制。模块的3.3V供电与系统共用同一路数字电源轨，上电完成之后面板上的指示灯呈点亮状态用以表明模块已进入正常工作阶段，此后主控随即通过UART接口向模块发出5字节初始化握手协议帧（$0xAA\ 0x55\ 0xFF\ 0x67\ 0xFB$）完成通信链路的建立。当需要触发TTS语音播报时，主控通过UART发送函数将对应的播报命令帧填入发送缓冲区；当模块侧完成语音识别并产生识别结果时，利用UART的RXNE接收中断逐字节收集完整的5字节应答帧并校验帧尾后提取出命令词的ID编码，供上层的应用逻辑以轮询方式读取。

全部命令词和播报词表按照固定格式在Excel表格中编辑完成后上传至云知声公司提供的语音识别固件及SDK开发平台，由该平台负责编译并输出语音播报固件文件，再由PC端工具通过有线方式将该固件烧录至语音模块的片内Flash当中以供运行使用。当前版本中，模块固件总计涵盖46条独立的TTS播报词以及10条可识别的命令词，这些条目完整覆盖了全部训练流程中的各类提示信息、反馈用语和控制指令。

语音交互集成模块如图\ref{fig:voice_module}所示。

\begin{figure}[h!]
\centering
\includegraphics[width=0.8\textwidth]{voice.jpg}
\caption{亚博智能AI语音交互集成模块}
\label{fig:voice_module}
\end{figure}

\section{按键与报警电路}

\subsection{患者反馈按键}

为满足训练过程中患者主动交互反馈的需求，系统在头戴装置外部设置了一颗患者反馈按键，连接MCU的PA2引脚。该按键采用轻触式微动开关，常开状态，按下时接通PA2到GND。PA2在CubeMX中配置为GPIO输入模式，内部上拉使能——未按下时引脚电平为逻辑高（上拉至3.3V），按下时引脚电平为逻辑低（接地）。

按键的去抖动和长短按检测由软件层面的调用按键扫描消抖函数实现。该函数在主循环中以10ms固定周期调用，读取PA2的当前电平并与前一帧状态比较：当检测到下降沿（高电平$\rightarrow$低电平）时记录按下时间戳；当检测到上升沿（低电平$\rightarrow$高电平）时，若按下的持续时间超过30ms的去抖动阈值，则判定为有效短按事件。

\subsection{蜂鸣器}

核心板板载了一枚无源压电蜂鸣器，由TIM12定时器的PWM输出驱动，连接于PB15引脚。占空比寄存器CCR2的值决定蜂鸣器的音量——CCR2=0时蜂鸣器静音，CCR2=500时输出50\%占空比的方波信号，此时蜂鸣器达到最大响度。

蜂鸣器报警通过控制报警函数的参数设置鸣叫方式，报警函数参数包括鸣叫次数、每次持续时间、间隔静音时长三个参数，通过循环交替设置和清零寄存器并结合延时生成特定的鸣叫模式。典型的报警方式为短促两声（各150ms，间隔100ms），与WS2812红色灯光和语音警告构成声-光-语音三维报警体系。

\section{电源与头戴式结构设计}

\subsection{电源系统}

整套装置经由XT30接口接入外部的直流电源进行供电，容许的输入电压覆盖了12V至24V的较宽范围。DM-02核心板在内部安排了两级DC-DC降压转换器来协同完成两级电压调节：前一级转换器负责将外部输入的高电压降至5V，在本设计中该路5V动力电源主要供给舵机、UART通信端口以及激光模组等需要较大驱动电流的外围设备使用；后一级转换器则在前一级5V的基础上进一步降压至3.3V，该路3.3V逻辑电源服务于STM32主控芯片、BMI088惯性传感器、CI1302语音模块以及其他工作在3.3V电压的数字电路。整机的电源树参见图\ref{fig:power_tree}所示。

\begin{figure}[h!]
\centering
\includegraphics[width=0.9\textwidth]{power.drawio.pdf}
\caption{电源树}
\label{fig:power_tree}
\end{figure}

需要特别指出的一点是，核心板的5V动力电源输出并不是上电后自动导通的，它的使能控制由三条独立的GPIO开关线路分别负责——PC13、PC14和PC15这三个引脚各自管理着一路电源通道的通断。这三路开关的极性统一设定为高电平有效，系统在完成对GPIO初始化函数的调用之后，随即将这三枚引脚统一配置为推挽输出模式并同时将其输出电平置为高状态，以保证各外围设备在启动初期就能够获得稳定的5V供电。

\subsection{头戴式结构设计}

由于本装置现阶段仍处于原理验证的原型机形态，考虑到结构设计环节的物料成本、开发时间周期与佩戴功能性之间的平衡需求，装置的承力外壳并未选择从空白开始建模并完整3D打印整个头戴式框架，而是直接借用了市面上现有的商业化自行车运动头盔作为基础承载平台。自行车运动头盔在材料选型和结构造型上本身就兼顾了轻量化、通风透气性以及一定的抗冲击防护能力，其壳体表面遍布多个通风孔洞，这些开孔位置恰好为线缆走线以及后续传感器和定位附件的扩展安装提供了天然的路径。成品头盔还自带了可以调节松紧程度的头围固定绑带与下颌束带，能够较好地适应不同佩戴个体头围尺寸的差异，为一次性原型验证提供了足够灵活的可调节佩戴舒适度。

各个子系统模块在头盔表面的空间布局大致遵从了以下几项原则：整体重心的分配尽量向头盔的中轴和后部靠拢以缓解单侧压迫的不适感；传感器的朝向严格与软件层面所定义的坐标系对齐以尽可能缩减标定阶段的偏差量；相连接模块之间的走线路径尽量压缩长度并加以绑扎以避免悬空晃荡对运动检测造成机械干扰。

开发板（DM-MC02核心板）的主安装位置被安排在了头盔的后部区域，与佩戴者后脑勺的位置相对应。固定姿态为芯片面朝外、USB接口竖直朝上。这种布局安排一方面是为了让板上BMI088芯片的X敏感轴自然而然地指向垂直方向（即USB口向上的方向），从而与第\ref{tab:axis_mapping}节中叙述的坐标系映射方案在物理层面保持直接的对应关系；另一方面也能够将开发板上分布较为集中的引出线缆统一收敛在头盔的后方，使装置正面区域在视觉上更为简洁干净。

装置中最核心的功能组件——SG90双轴舵机云台以及激光模组——借助云台底座牢牢固定在头盔的顶部区域。云台底座的中央位置预留了一个与市面通行的GoPro运动相机头盔绑带式支架完全兼容的1/4英寸标准螺孔，这一接口规格使得整个云台模块可以通过绑带紧固的纯机械方式牢靠地连接在头盔顶部，而不需要在头盔壳体上进行打孔或者使用胶粘剂进行不可逆的固定操作。该连接方式不仅提供了足够的机械约束力来确保训练过程中的稳定指向，还赋予了云台组件在需要时可以快速拆卸与重新安装的灵活度，为日常调试和后续维护操作提供了较大的便捷性。

语音交互模块被安放在了头盔的前部区域，位置大致对应佩戴者的前额附近。这种布局有助于缩短模块自带的微型麦克风与佩戴者口腔之间的直线距离，从而对语音指令的识别灵敏度产生积极的提升作用，同时也使得模块内置的微型扬声器与佩戴者耳部之间保持了适中的间距，既保证了语音播报内容的听觉清晰度，又避免了因音量过大而对使用者造成听觉疲劳。

患者反馈按键与可选的备用电池组共同集成于一只紧凑的小型独立控制盒的内部，该控制盒通过连接线与固定于头盔后方的开发板之间建立电气连接。控制盒本身同样使用PLA材料经3D打印工艺制作成形，其内部空间被设计得较为紧凑，方便患者单手握持并进行按键操作。在当前的原型验证阶段，系统以有线方式接入外部直流电源作为主要的供电手段，备选电池选用了一颗市面上通用的航模规格锂电池，其详细的技术参数汇总于表\ref{tab:battery_parameters}中。

\begin{table}[h!]
\centering
\caption{航模电池参数}
\label{tab:battery_parameters}
\begin{tabular}{p{4cm} p{4cm}}
\toprule[1.5pt]
\textbf{特性} & \textbf{参数} \\
\midrule[0.75pt]
容量&1500mAh\\
电压&14.8V\\
电芯&4s1p\\
插头&XT30\\
重量&155g±10g\\
尺寸&68*30*34mm\\
\bottomrule[1.5pt]
\end{tabular}
\end{table}

聚焦训练装置采用一套折叠连杆机构，安装于头盔前部正中且靠近佩戴者鼻梁上方的位置，专门用于第4章所述的视觉聚焦训练模式（模式D）。该装置的核心是一个红色LED灯珠，被固定于连杆机构的最前端。连杆机构由三节长约5cm的轻质杆件通过微型铰链顺序串接而成，其末端的旋转基座固定于头盔前部的一个支座上。连杆在非使用状态下可以完全折叠并贴附于头盔前部的轮廓曲面之上；进入聚焦训练时，患者可用单手将连杆逐步展开，使得LED灯珠移动至距离眼睛约15cm的近焦注视位置。由于铰链与基座的配合面带有适量的阻尼，连杆可在行程范围之内的任意角度实现停驻自锁，便于根据患者不同视力水平对近焦距离进行精细微调。装置采用PLA材料经3D打印工艺制成，全部构件的累计重量不足20g，日常的展开与折叠均可用单手独立完成。

除了借用的成品自行车头盔本体之外，装置中所有需要量身定制的结构零部件——包括开发板固定支架、云台底座、控制盒外壳以及聚焦训练装置的整套连杆机构——均借助拓竹品牌的3D打印机使用PLA（聚乳酸）材料逐层打印而成。3D打印工艺的单件成型周期较为短暂、原材料的单位成本也相对低廉，因此在原型验证阶段有利于实现快速的设计迭代和及时的尺寸修正。

装置包含核心板、云台组件、激光模组、语音模块、成品头盔以及全部连接线缆在内的累计总重量经过实际称量后约为500g，处于人体头部可长期佩戴而无显著不适的重量范围之内。
